\expandafter\def\csname HistoricalNoteText16\endcsname{[Handschriftliche Bemerkung.] Am 31. Dec. schon fast 11 Grad vergl. p. 581 u. 647. [S. 203 d. B.]}
\clearpage\SourceLeaf{204}%
\pdfbookmark[1]{Fortgesetzte Nachrichten. Ceres. 1801 December}{ceres180112}
\fancyhead[L]{\small [CERES.]}
\begin{center}
{\small FORTGESETZTE NACHRICHTEN ÜBER DEN LÄNGST VERMUTHETEN NEUEN HAUPT-PLANETEN UNSERES SONNEN-SYSTEMS.\par}
[CERES.]\par
\medskip\hrule\medskip
{\small Monatliche Correspondenz zur Beförderung der Erd- und Himmels-Kunde, herausgegeben vom Freiherrn \textsc{von Zach}. Band III. 1801 December.\par}
\medskip\hrule
\end{center}\medskip
\noindent —\quad —\quad —\par
Eine grosse Hoffnung der Hülfe und Erleichterung gewähren uns die uns jüngst mitgetheilten Untersuchungen und Berechnungen des Dr.\ \textsc{Gauss} in Braunschweig. Sie geben uns zugleich einen neuen und hohen Grad der Wahrscheinlichkeit, dass das von \textsc{Piazzi} entdeckte neue Gestirn ein wirklich planetarischer Weltkörper sei, welcher sich nach \textsc{Kepler}'schen Gesetzen zwischen der Mars- und Jupiters-Bahn fortbewegt.

Wir eilen mit der Mittheilung seiner Berechnungen um so mehr, da seine neue elliptische Bahn von der elliptischen des Dr.\ \textsc{Burckhardt}, und von den zwei Kreis-Bahnen des Dr.\ \textsc{Olbers} und P.\ \textsc{Piazzi}, welche wir in unseren vorigen Heften, mit hieraus im voraus berechneten Orten des Planeten, mitgetheilt haben, beträchtlich verschieden ist, und ihre Abweichung, in gegenwärtigem Monat, von den \textsc{Gauss}'schen Positionen sich auf 6 bis 7 Grade in geocentrischer Länge erstrecken kann. Es ist daher von Wichtigkeit, diese Bemerkungen den practischen Astronomen bald möglichst mitzutheilen, weil sie dadurch erfahren, dass sie \textit{nothwendig} den Raum am Himmel, worin sie dieses neue und so schwer zu findende Gestirn aufzusuchen haben, um 6 bis 7 Grade weiter nach Osten ausdehnen müssen. Denn 1) liegen nach Dr.\ \textsc{Gauss}'s Rechnung die \textsc{Piazzi}'schen Beobachtungen, wie Dr.\ \textsc{Burckhardt} und \textsc{Olbers} vorausgesetzt haben, weder nahe am Perihelium, noch nahe beim Aphelium, sondern fast mitten zwischen beiden.

2) Ist die Excentricität der Bahn nach seiner Rechnung nicht so ganz unbeträchtlich, wie P.\ \textsc{Piazzi} glaubt, daher denn die nach seinen elliptischen Elementen im voraus berechneten Längen von den nach Kreis-Hypothesen berechneten in gegenwärtigem Monat leicht um 7 Grade\HistoricalNote{16} verschieden sein können.

3) Rechtfertigen die \textsc{Gauss}'schen Elemente die \textsc{Piazzi}'schen Beobachtungen vollkommen, und befreien sie ganz von dem Verdachte eines Mangels an Genauigkeit, wie Dr.\ \textsc{Burckhardt} und \textsc{Olbers} bei der Prüfung durch Differenzen vermuthet hatten.

Alles dieses beweist die \textsc{Gauss}'sche Ellipse. Welches Vertrauen sie daher erwecken muss, werden Astronomen aus der Genauigkeit erkennen, mit welcher sie die sämmtlichen \textsc{Piazzi}'schen Beobachtungen darstellt. Dr.\ \textsc{Gauss} ist zu diesen Berechnungen durch einige Untersuchungen über physische Astronomie
\SourceLeaf{205}%
veranlasst worden, welche ihn auf einige nicht unerhebliche Zusätze zur Theorie der Bestimmung der Himmelskörper in Kegelschnitten jeder Art geführt hat, und wovon er uns einiges mitzutheilen so gütig war. Wir hoffen unsere astronomischen Leser ein andermal damit zu unterhalten, da uns diese Auseinandersetzung hier zu weit von unserm Gegenstande ablenken würde; wir schränken uns daher blos auf das ein, was unmittelbar auf die Untersuchung der Bahn der \textit{Ceres Ferdinandea} Bezug hat.

Zuerst wählte Dr.\ \textsc{Gauss} zur Bestimmung der Bahn die drei Beobachtungen vom 2. und 22. Jan. und vom 11. Febr., wobei er diese Data ganz so angenommen hat, wie sie in dem Septbr.\ St.\ d.\ M.\ C.\ S.\ 280 gegeben sind. Nach einem eigenthümlichen Verfahren fand er sofort beim vierten Versuche folgende Elemente:
\begin{center}
{\small\setlength{\tabcolsep}{2pt}
\begin{tabular}{l r}
Sonnenferne $\ .\ .\ .\ .\ .\ .\ .\ .\ .\ .$&$330^0\ 14'\ 33''$\\
$\AscendingNode\ .\ .\ .\ .\ .\ .\ .\ .\ .\ .\ .\ .\ .$&$81\quad8\ 50$\\
Neigung der Bahn $\ .\ .\ .\ .\ .\ .\ .\ .$&$10\ 32\ 19$\\
Logarithmus der halben grossen Axe $\ .\ .$&$0.4381058$\\
Excentricität $\ .\ .\ .\ .\ .\ .\ .\ .\ .\ .$&$0.0832836$\\
\multicolumn{2}{l}{Epoche: 31. Dec. 1800 im mittl.\ Mittag zu Palermo}\\
mittlere heliocentrische Länge: $\ .\ .\ .$&$77^0\ 54'\ 29''$\\[4pt]
\multicolumn{2}{c}{Hieraus folgt:}\\[3pt]
grösste Mittelpunkts-Gleichung $\ .\ .\ .$&$9^0\ 32'\ 57''$\\
halbe grosse Axe $\ .\ .\ .\ .\ .\ .\ .\ .$&$2.74226$\\
siderische Umlaufszeit $\ .\ .\ .\ .\ .\ .$&$1658\frac23$ Tage\\
tägliche siderische mittlere Bewegung&$781''.355$
\end{tabular}}
\end{center}
Diese Elemente stellen die beiden äussern Beobachtungen genau, die mittlere mit $2''$ Fehler in Länge und Breite dar. Auf Aberration und Präcession ist gehörig Rücksicht genommen worden.

Nach diesem ersten glücklichen Versuche unternahm Dr.\ \textsc{Gauss} eine zweite Berechnung dieser Bahn. Er liess die Beobachtungen vom 2. und 22. Jan. weg, und wählte anstatt derselben die vom 1. und 21. Jan., in Verbindung mit der vom 11. Febr., welche er nicht weglassen wollte, damit die Zwischenzeit so gross als möglich bliebe. Diese Rechnung hat ihm, bei der vierten Hypothese, folgende Elemente gegeben:
\begin{center}
{\small\setlength{\tabcolsep}{2pt}
\begin{tabular}{l r}
Sonnenferne $\ .\ .\ .\ .\ .\ .\ .\ .\ .$&$330^0\ 33'\ 20''$\\
$\AscendingNode\ .\ .\ .\ .\ .\ .\ .\ .\ .\ .\ .\ .\ .$&$81\quad2\ 35$\\
Neigung der Bahn $\ .\ .\ .\ .\ .\ .\ .$&$10\ 36\ 30$\\
Logarithmus der halben grossen Axe&$0.4370335$\\
Excentricität $\ .\ .\ .\ .\ .\ .\ .\ .\ .$&$0.0705553$\\
Epoche wie oben $\ .\ .\ .\ .\ .\ .\ .$&$76^0\ 28'\ 14''.27$\\[4pt]
\multicolumn{2}{c}{Hieraus folgt:}\\[3pt]
grösste Mittelpunkts-Gleichung $\ .\ .\ .$&$8^0\ 5'\ 19''$\\
halbe grosse Axe $\ .\ .\ .\ .\ .\ .\ .\ .$&$2.73548$\\
siderische Umlaufszeit $\ .\ .\ .\ .\ .\ .$&$1652\frac12$ Tage\\
tägliche siderische mittlere Bewegung $\ .$&$784''.25$
\end{tabular}}
\end{center}
Nach diesen Elementen stimmen sämmtliche im September-Hefte S.\ 280 angeführte \textsc{Piazzi}'sche Beobachtungen folgendermaassen:
\begin{center}
{\small\setlength{\tabcolsep}{4pt}
\begin{tabular}{r|r|r|r|r}
1801&\multicolumn{2}{c|}{Berechnete}&\multicolumn{2}{c}{Fehler der}\\
&Länge&Breite&Länge&Breite\\\hline
Jan.\ 1&$53^0\ 22'\ 58''42$&$3^0\ 6'\ 42''09$&$+0''12$&$-0''01$\\
2&$53\ 19\ 37.02$&$3\ 2\ 23.78$&$-7.28$&$-1.12$\\
$\mathord{:}\mathord{:}$ 3&$53\ 16\ 43.67$&$2\ 58\ 6.70$&$-14.93$&$-3.20$\\
4&$53\ 14\ 14.03$&$2\ 53\ 51.16$&$-1.47$&$-4.44$\\
$\mathord{:}\mathord{:}$ 10&$53\ 7\ 54.51$&$2\ 28\ 53.13$&$-4.59$&$-7.47$\\
13&$53\ 10\ 18.77$&$2\ 16\ 48.78$&$-18.83$&$-10.92$\\
14&$53\ 11\ 55.25$&$2\ 12\ 51.23$&$-5.95$&$-5.47$\\
19&$53\ 26\ 0.37$&$1\ 53\ 34.19$&$+1.17$&$-4.01$\\
21&$53\ 34\ 22.68$&$1\ 46\ 5.98$&$+1.28$&$-0.02$\\
22&$53\ 39\ 7.88$&$1\ 42\ 25.06$&$+6.08$&$-3.04$
\end{tabular}
\par\medskip
\begin{tabular}{r|r|r|r|r}
1801&\multicolumn{2}{c|}{Berechnete}&\multicolumn{2}{c}{Fehler der}\\
&Länge&Breite&Länge&Breite\\\hline
Jan.\ 23&$53^0\ 44'\ 15''74$&$1^0\ 38'\ 46''25$&$+0''04$&$-5''85$\\
28&$54\ 15\ 20.88$&$1\ 21\ 4.07$&$+5.18$&$-2.83$\\
30&$54\ 30\ 4.25$&$1\ 14\ 13.73$&$+5.25$&$-2.27$\\
31&$54\ 38\ 11.25$&$1\ 10\ 51.99$&$+3.95$&$-2.61$\\
$\mathord{:}\mathord{:}$ Feb.\ 1&$54\ 46\ 28.37$&$1\ 7\ 32.09$&$+9.07$&$+1.19$\\
2&$54\ 55\ 5.09$&$1\ 4\ 14.25$&$+7.19$&$+3.75$\\
5&$55\ 22\ 50.25$&$0\ 54\ 32.88$&$+6.85$&$+3.98$\\
8&$55\ 53\ 23.15$&$0\ 45\ 9.20$&$-6.38$&$+4.20$\\
11&$56\ 26\ 39.97$&$0\ 36\ 2.90$&$-0.03$&$+0.00$
\end{tabular}}
\end{center}

\SourceLeaf{206}%
So vortrefflich diese Übereinstimmung ist, so hält es Dr.\ \textsc{Gauss} doch selbst für sehr möglich, dass seine Elemente sich von den wahren beträchtlich mehr entfernen können, als sie unter sich abweichen, da der durchlaufene Theil der Bahn so klein ist, und von der ersten bis zur letzten Beobachtung nur $9^0\ 15'\ 35''$ beträgt. Indessen da diese Bahn doch \textit{möglich}, wo nicht \textit{höchst wahrscheinlich} ist: so haben practische Astronomen doppelt Ursache, auf dieselbe bei Aufsuchung dieses Gestirns Rücksicht zu nehmen und sie zu ihrem vorzüglichsten Leitfaden zu wählen, da von allen bisher versuchten und berechneten Bahnen keine eine solche nahe und genaue Vereinigung aller Beobachtungen darstellt, als eben diese \textsc{Gauss}'sche Ellipse.

Übrigens bemerken wir noch, dass Dr.\ \textsc{Gauss}, bei allen diesen Berechnungen sowol auf den im October-Hefte S.\ 365 angezeigten Druckfehler von einer Minuta in dem Orte der Sonne, als auch auf einen zweiten beim 13. Jan., ebenfalls bei der Länge der Sonne vorgefallenen, Rücksicht genommen, und diese letzte zu $9^{\mathrm z}\ 23^0\ 13'\ 13.''8$ vorausgesetzt habe.

So weit war Dr.\ \textsc{Gauss} in seinen Berechnungen der Bahn der \textit{Ceres Ferdinandea} gekommen, als er sie uns mitzutheilen so gütig war. Wir hatten indessen die \textsc{Piazzi}'sche Abhandlung, nebst seinen verbesserten Beobachtungen erhalten, wovon wir im November-Hefte einen getreuen Auszug mitgetheilt haben. So gering aber auch diese Veränderungen, worunter die um $15''$ verminderte gerade Aufsteigung vom 11. Febr. die beträchtlichste ist, und wodurch die vorigen Elemente nur wenig geändert werden: so haben sich doch hier und da einige bedeutendere Druck- und Rechnungsfehler bei der Reduction dieser Beobachtungen eingeschlichen; und da sich diese Fehler leicht wegschaffen, und die Elemente hiernach bald verbessern liessen: so theilten wir dem Dr.\ \textsc{Gauss} die sämmtlichen \textsc{Piazzi}'schen Beobachtungen in der Italienischen Originalschrift zu diesem Behufe mit. Allein ehe er noch unsere Zuschrift erhielt, hatte er bereits den Anfang mit der Berechnung einer neuen Bahn gemacht. So gering auch die Abweichungen seiner zweiten Elemente von den Beobachtungen schon waren: so hatten sie doch noch eine sehr auffallende Regularität, woraus sich mit Gewissheit vorhersehen liess, dass man die Fehler noch beträchtlicher vermindern könne. Er machte den Versuch, und da er beim Empfang der ihm zugeschickten verbesserten Beobachtungen, so wie sie \textsc{Piazzi} selbst reducirt und berechnet hatte, diese Arbeit schon grösstentheils zu Ende gebracht hatte: so hielt er es für gut, sie erst ganz zu vollenden, um so mehr, da er voraus wusste, dass die Länge nach den neuen Elementen für den 11. Febr. um $6''$ kleiner ausfallen, und sich also der berichtigten, ihm bis dahin natürlich unbekannten Angabe der Beobachtung schon von selbst fast um die Hälfte nähern würde. Auf diese Art fand er folgende \textit{dritte} Elemente:
\begin{center}
{\small\setlength{\tabcolsep}{2pt}
\begin{tabular}{l r}
Sonnenferne $\ .\ .\ .\ .\ .\ .\ .\ .\ .\ .$&$326^0\ 53'\ 50''$\\
$\AscendingNode\ .\ .\ .\ .\ .\ .\ .\ .\ .\ .\ .\ .\ .$&$81\quad1\ 44$\\
Neigung der Bahn $\ .\ .\ .\ .\ .\ .\ .\ .$&$10\ 36\ 21$\\
Logarithmus der halben grossen Axe $\ .$&$0.4414902$\\
Excentricität $\ .\ .\ .\ .\ .\ .\ .\ .\ .\ .$&$0.0819603$\\
Epoche $\ .\ .\ .\ .\ .\ .\ .\ .\ .\ .\ .\ .$&$77^0\ 34'\ 28''$\\[4pt]
\multicolumn{2}{c}{Hieraus ferner:}\\[3pt]
Grösste Mittelpunkts-Gleichung $\ .\ .\ .$&$9^0\ 23'\ 57''$\\
Tägliche siderische mittl.\ hel.\ Bewegung&$772''275$\\
Tägliche tropische helioc.\ Bewegung $\ .$&$772.413$\\
Tropische Umlaufszeit $\ .\ .\ .\ .\ .\ .$&$1677.8$ Tage
\end{tabular}}
\end{center}
Diese Elemente stellen die von uns reducirten und im September-Hefte S.\ 280 angeführten \textsc{Piazzi}'schen Beobachtungen folgendermaassen dar:
\expandafter\def\csname HistoricalNoteText17\endcsname{In den \textsc{Piazzi}'schen Original-Beobachtungen S.\ 565 des November-Stücks sind die Längen der Sonne am 13. Januar und 8. Februar um eine Minute zu gross; alle geocentrische Breiten gegen $10''$ zu klein, welches ohne Zweifel daher kommt, dass \textsc{Piazzi} bei der Reduction seiner Beobachtungen sich einer andern Schiefe der Ekliptik und wahrscheinlich jener aus der \textit{Conn.\ d.\ tems Année} IX. zu $23^0\ 28'\ 1''$ bedient hat.}
\SourceLeaf{207}%
\begin{center}
{\small\setlength{\tabcolsep}{4pt}
\begin{tabular}{r|r|r|r|r}
1801&\multicolumn{2}{c|}{Berechnete}&\multicolumn{2}{c}{Fehler der}\\
&Länge&Breite&Länge&Breite\\\hline
Jan.\ 1&$53^0\ 23'\ 2''34$&$3^0\ 6'\ 43''63$&$+4''04$&$+1''53$\\
2&$53\ 19\ 41.24$&$3\ 2\ 25.68$&$-3.06$&$+0.78$\\
3&$53\ 16\ 48.05$&$2\ 58\ 8.97$&$-10.35$&$-0.93$\\
4&$53\ 14\ 18.47$&$2\ 53\ 53.79$&$+2.97$&$-1.81$\\
10&$53\ 7\ 58.37$&$2\ 28\ 57.12$&$-0.73$&$-3.48$\\
13&$53\ 10\ 21.60$&$2\ 16\ 52.89$&$-16.00$&$-6.81$\\
14&$53\ 11\ 57.70$&$2\ 12\ 55.36$&$-3.50$&$-1.34$\\
19&$53\ 26\ 0.59$&$1\ 53\ 38.01$&$+1.39$&$-0.19$\\
21&$53\ 34\ 21.99$&$1\ 46\ 9.53$&$+0.69$&$+3.53$\\
22&$53\ 39\ 6.69$&$1\ 42\ 28.45$&$+4.89$&$+0.35$
\end{tabular}
\par\medskip
\begin{tabular}{r|r|r|r|r}
1801&\multicolumn{2}{c|}{Berechnete}&\multicolumn{2}{c}{Fehler der}\\
&Länge&Breite&Länge&Breite\\\hline
Jan.\ 23&$53^0\ 44'\ 14''08$&$1^0\ 38'\ 49''44$&$-1''62$&$-2''66$\\
28&$54\ 15\ 17.11$&$1\ 21\ 5.91$&$+1.41$&$-0.99$\\
30&$54\ 30\ 9.76$&$1\ 14\ 15.12$&$+0.76$&$-0.88$\\
31&$54\ 38\ 6.44$&$1\ 10\ 52.81$&$-0.86$&$-1.79$\\
Febr.\ 1&$54\ 46\ 23.22$&$1\ 7\ 32.54$&$+3.92$&$+1.64$\\
2&$54\ 54\ 59.71$&$1\ 4\ 14.30$&$+1.81$&$+3.80$\\
5&$55\ 22\ 44.30$&$0\ 54\ 31.72$&$+0.90$&$+2.82$\\
8&$55\ 53\ 17.01$&$0\ 45\ 6.65$&$-12.49$&$+1.63$\\
11&$56\ 26\ 34.10$&$0\ 35\ 58.96$&$-5.90$&$-3.94$
\end{tabular}}
\end{center}
Die von \textsc{Piazzi} angezeigten Fehler und Verbesserungen geben nunmehr folgende Abänderungen: Wegen der um $15''$ verminderten geraden Aufsteigung vom 11. Febr. wird, mit $23^0\ 28'\ 12''$ Schiefe der Ekliptik, die Länge $56^0\ 26'\ 26''1$, die Breite $35'\ 59''7$: also Fehler der Länge $+8''0$, Fehler der Breite $-0''74$. Beim Nachrechnen der Reduction der geraden Aufsteigung und Abweichung bei denjenigen Beobachtungen, wo \textsc{Piazzi}'s Angaben beträchtlich von den unsrigen abweichen, fand Dr.\ \textsc{Gauss} die Länge am 8. Febr. $55^0\ 53'\ 17''7$: folglich ist der Fehler $-0''7$\HistoricalNote{17}. Das übrige, und auch selbst die bei \textsc{Piazzi} so stark abweichende Länge am 13. Januar findet Dr.\ \textsc{Gauss}, wie bei uns. Da diese Beobachtung sowohl von diesen letzten als auch von den vorigen Elementen gerade am meisten abweicht und sich nicht genau darstellen lässt, ohne den übrigen Beobachtungen Zwang anzuthun: so vermuthet Dr.\ \textsc{Gauss}, dass irgend dabei ein kleines Versehen vorgefallen sei. Übrigens sieht man leicht, dass die übrigen von \textsc{Piazzi} angegebenen kleinen Verbesserungen (Novembr.-St.\ S.\ 573) die Genauigkeit, womit die Beobachtungen durch diese letzten Elemente dargestellt werden, wenig ändert.

So gering indessen alle Fehler sind, so hat Dr.\ \textsc{Gauss} doch noch einen vierten Versuch gemacht, die Genauigkeit noch etwas zu vergrössern. Die Längenfehler sind bei diesen letzten Elementen im Februar sämmtlich positiv, wenn man die kleinen Verbesserungen von \textsc{Piazzi} noch in Betrachtung zieht; auch die Breitenfehler neigen sich gegen das Ende so ziemlich nach einer Seite. Dr.\ \textsc{Gauss} hat daher neue Elemente gesucht, wonach die berechneten Längen und Breiten im Februar etwas kleiner werden, indem die übrigen ungefähr dieselbe Grösse behalten. Ob er gleich diese Elemente mit den sämmtlichen Beobachtungen nicht verglichen hat, welches zum Theil auch ganz überflüssig ist: so glaubt er doch im voraus versichern zu können, dass sie zwischen den Beobachtungen, so viel sich thun lässt, das Mittel halten, und dass die Fehler bei keiner Beobachtung (die vom 3. und 13. Januar ausgenommen) über $5''$ in Länge und Breite gehen, gar keine Regularität mehr haben, sondern eine ganz irreguläre Abwechselung der Zeichen geben. Dr.\ \textsc{Gauss} zweifelt nicht ohne Grund, ob man \textit{bei eben den Datis} eine \textit{merklich grössere} Übereinstimmung bei andern Elementen finden könne, womit er indessen durchaus nicht sagen oder verstanden wissen will, dass andere erheblich verschiedene Elemente eine ähnliche Übereinstimmung nicht gewähren könnten; zumal wenn bei dieser delicaten Rechnung, wo wenige Secunden schon einen merklichen Ausschlag geben, etwas andere Bestimmungen von der Sonne gebraucht würden. So ist die Länge der Sonne bei \textsc{Piazzi} im Februar von der unsrigen um eine halbe Minute verschieden, wiewohl gerade hier in der Nähe der Quadratur der Einfluss der Sonnen-Länge geringer ist, als in andern Lagen. Dr.\ \textsc{Gauss} glaubt daher, dass es nicht undienlich wäre, wenn man die Fehler der Sonnen%
\SourceLeaf{208}%
tafeln aus sehr genauen Beobachtungen für diese Zeiten bestimmmte, und die Örter der Sonne hienach verbesserte. Diese \textit{vierten} Elemente sind nun folgende:
\begin{center}
{\small\setlength{\tabcolsep}{2pt}
\begin{tabular}{l r}
Sonnenferne $\ .\ .\ .\ .\ .\ .\ .\ .\ .\ .$&$326^0\ 27'\ 38''$\\
$\AscendingNode\ .\ .\ .\ .\ .\ .\ .\ .\ .\ .\ .\ .\ .$&$81\quad0\ 44$\\
Neigung $\ .\ .\ .\ .\ .\ .\ .\ .\ .\ .\ .$&$10\ 36\ 57$\\
Logarithmus der halben grossen Axe $\ .$&$0.4420527$\\
Excentricität $\ .\ .\ .\ .\ .\ .\ .\ .\ .\ .$&$0.0825017$\\
Epoche 1800 31. Dec.\ $\ .\ .\ .\ .\ .\ .\ .$&$77^0\ 36'\ 34''$\\[4pt]
\multicolumn{2}{c}{Hieraus:}\\[3pt]
grösste Mittelpunkts-Gleichung $\ .\ .\ .\ .$&$9^0\ 27'\ 41''$\\
\multicolumn{2}{l}{tägliche mittlere helioc.\ tropische Bewegung}\\
$\ .\ .\ .\ .\ .\ .\ .\ .\ .\ .\ .\ .$&$770''914$
\end{tabular}}
\end{center}
Aus diesen Elementen hat Dr.\ \textsc{Gauss} folgende Örter der \textit{Ceres Ferdinandea} im voraus berechnet. Die Zeit ist mittlere für Mitternacht in \textit{Palermo}.
\begin{center}
{\fontsize{9}{11}\selectfont\setlength{\tabcolsep}{3pt}
\begin{tabular}{r|r|r|r|r|r}
1801&\shortstack{Geocentri-\\sche Länge}&\shortstack{Geocentri-\\sche Breite\\nördlich}&\shortstack{Logarithm.\\des Ab-\\standes von\\der $\EarthSymbol$}&\shortstack{Logarithm.\\des Ab-\\standes von\\der $\odot$}&\shortstack{Verhält-\\niss der ge-\\sehenen\\Helligkeit}\\\hline
November 25&$5^{\mathrm z}\ 20^0\ 16'$&$9^0\ 25'$&$0.42181$&$0.40468$&$0.6102$\\
December 1&$5\ 22\ 15$&$9\ 48$&$0.40940$&$0.40472$&$0.6459$\\
7&$5\ 24\ 7$&$10\ 12$&$0.39643$&$0.40479$&$0.6855$\\
13&$5\ 25\ 51$&$10\ 37$&$0.38296$&$0.40488$&$0.7290$\\
19&$5\ 27\ 27$&$11\ 4$&$0.36902$&$0.40499$&$0.7770$\\
25&$5\ 28\ 53$&$11\ 32$&$0.35468$&$0.40512$&$0.8295$\\
31&$6\ 0\ 10$&$12\ 1$&$0.34000$&$0.40528$&$0.8869$
\end{tabular}}
\end{center}
Sollte man den Ort des Planeten nach diesen Elementen genauer, oder auf eine längere Zeit berechnen wollen: so setzen wir zu diesem Behufe noch folgende Formeln hierher:

\begin{center}1) Zur Berechnung der Mittelpunkts-Gleichung $=M$\end{center}
\[
\begin{aligned}
M={}&-34005''494\ \mathrm{Sin.\ Anom.\ med.}\\
&+1750''951\ \mathrm{Sin.}\ 2\ \mathrm{A.\ m.}-124''951\ \mathrm{Sin.}\ 3\ \mathrm{A.\ m}\\
&+10''192\ \mathrm{Sin.}\ 4\ \mathrm{A.\ m.}-0''901\ \mathrm{Sin.}\ 5\ \mathrm{A.\ m.}\\
&+0''083\ \mathrm{Sin.}\ 6\ \mathrm{A.\ m.}
\end{aligned}
\]
\begin{center}2) Für den Radius Vector $=r$\end{center}
\[
\begin{aligned}
r={}&2.776695+0.2373586\ \mathrm{Cosin.\ Anom.\ med.}\\
&-0.0093752\ \mathrm{Cosin.}\ 2\ \mathrm{A.\ m.}\\
&+0.0005789\ \mathrm{Cosin.}\ 3\ \mathrm{A.\ m.}\\
&-0.0000423\ \mathrm{Cosin.}\ 4\ \mathrm{A.\ m.}\\
&+0.0000034\ \mathrm{Cosin.}\ 5\ \mathrm{A.\ m.}\\
&-0.0000003\ \mathrm{Cosin.}\ 6\ \mathrm{A.\ m.}
\end{aligned}
\]
\[
\text{Oder:}\quad r=\frac{7.60570}{2.767278\pm0.2283053\ \mathrm{Cosin.\ Anom.\ ver.}}
\]
\begin{center}3) Für die heliocentrische Breite $=\lambda$.\end{center}
\[
\mathrm{Log.\ Sin.}\ \lambda=9.2653438+\mathrm{Log.\ Sin.\ Arg.\ Latit.}
\]
\begin{center}4) Für die Reduction der heliocentr.\ Länge auf die Erdbahn $=\varepsilon$.\end{center}
\[
\begin{aligned}
\alpha)&\quad\mathrm{Log.\ Tang.}\ \varphi=9.9925025+\mathrm{Log.\ Tang.\ Arg.\ Latit.}\\
\beta)&\quad\varepsilon=\mathrm{Arg.\ Latit.}-\varphi.
\end{aligned}
\]
\begin{center}5) Für die Reduction des Radius Vector, oder curtirte Distanz $=\rho$.\end{center}
\[
\rho=\mathrm{Cosin.\ Latit.\ helioc.}\times\mathrm{Distant.\ vera.}
\]
\begin{center}Oder für den Logarithm.\ der Verkürzung selbst $=\mathrm{Log.}\ \mu$.\end{center}
\[
\mathrm{Log.}\ \mu=10.000000-\mathrm{Log.\ Cosin.\ Latit.\ helioc.}
\]
\begin{center}6) Für die Aberration des Lichts, in Länge, Breite, Grade Aufsteigung und Abweichung $=a$.\end{center}
\[
\begin{aligned}
\mathrm{Log.}\ a={}&\pm\mathrm{Log.\ Dist.}\ \text{à}\ \EarthSymbol+\mathrm{Log.\ mot.\ horar.\ geoc.}\\
&+\mathrm{Log.}\ 7.751007.
\end{aligned}
\]

\SourceLeaf{209}%
Wir beschliessen diese Nachricht mit einer Bemerkung des Dr.\ \textsc{Gauss} über die Neigung der Bahn der \textit{Ceres}, welche wegen ihrer Grösse vielen Astronomen so auffallend war. Ob er gleich mit dem Urtheil einiger Astronomen ganz einverstanden ist, dass wir durch keinen physischen Grund berechtigt werden, bei allen Weltkörpern unseres Sonnen-Systems, die eine kreisähnliche Bahn haben, auch eine geringe Neigung gegen die Ekliptik zu erwarten: so scheint ihm doch, dass das Paradoxe davon beträchtlich vermindert werde, und die Analogie sich doch einigermassen retten lasse, wenn man nur diese Planeten-Bahn, wie \textsc{La Place} bei den Bahnen der Uranus-Trabanten gethan hat (\textit{A.\ G.\ E.}\ II B.\ S.\ 259) auf ihre eigentliche Ebene, auf welche man sie beziehen muss, nämlich auf den Sonnen-Aequator, bezogen haben wird. Wenn wir die Ebenen der sämmtlichen Planeten unter einander vergleichen, so zeigt sich alsdann, dass die Bahn der \textit{Ceres} gegen die Bahn keines einzigen andern Planeten so gross ist, als gegen die Erdbahn. Vergleichen wir also die Bahn der \textit{Ceres} mit der Erdbahn, so vergleichen wir gerade die beiden Extreme unseres Sonnen-Systems. Denkt man sich aber eine Ebene, die zwischen den Ebenen sämmtlicher \textit{acht} Planeten-Bahnen ungefähr in der Mitte liegt: so sind die Neigungen gegen dieselbe alle klein genug. Es scheint sehr merkwürdig, dass gerade der Sonnen-Aequator ungefähr eine solche Lage hat; nur mit der Einschränkung, dass die Bahn der \textit{Ceres} nebst denen der Venus und des Mercur am wenigsten, etwa 3 bis 4 Grade, die Bahn der Erde hingegen gerade am meisten dagegen geneigt ist.
\begin{flushright}[\textsc{von Zach}.]\end{flushright}

\bigskip\Needspace{10\baselineskip}
\pdfbookmark[1]{Fortgesetzte Nachricht. Ceres. 1802 Februar}{ceres180202}
\begin{center}
{\small FORTGESETZTE NACHRICHT ÜBER DEN NEUEN HAUPT-PLANETEN.\par}
[CERES.]\par
\medskip\hrule\medskip
Monatliche Correspondenz. 1802 Februar.\par
\medskip\hrule
\end{center}\medskip
—\quad —\quad —\quad „In der That bewundernswürdig ist die Übereinstimmung der Abweichungen der \textsc{Piazzi}'schen Beobachtungen, sowol unter sich, als mit der Rechnung, und gerade dies bestärkt meine Hoffnung, dass meine Elemente zur Auffindung hinreichend genau sein werden. Denn gerade die Breiten-Beobachtungen, die von den Abweichungen am meisten abhängen, haben auf die Zustimmung der Elemente den grössten Einfluss; eine Änderung von $20''$ in der Breite, wenn man sie bei den äussern Beobachtungen positiv und bei den mittleren negativ anbrächte, oder umgekehrt, würde die Ellipse total ändern. Mit dieser Bemerkung hängt eine andere zusammen, die aber dann erst ihren vollen Werth erhält, wenn die Wiederauffindung geglückt sein, und die aus den Beobachtungen abgeleiteten Elemente bestätigen wird, dass wir dieselbe in diesem Falle hauptsächlich der starken Neigung der Bahn verdanken werden. Fiele die Bahn mit der Ekliptik zusammen: so würde man wol darauf Verzicht thun müssen, wenigstens würde die Ungewissheit nach einem Jahr in der Länge ausserordentlich gross, und vielleicht ein Paar Zeichen betragen, obwol diese Zone alsdann fast gar keine Breite haben würde. Man sieht den Grund davon leicht ein, wenn man erwägt, dass alsdann zur Bestimmung der Ellipse nothwendig vier Längen erfordert würden, dass folglich zweien äussern und einer mittleren Beobachtung durch unendlich viele Ellipsen, und vielleicht auch Parabeln und Hyperbeln vollkommen Genüge geleistet werden könnte, und da bei der Kleinheit der Reihe der Beobachtungen leicht zu ermessen
\SourceLeaf{210}%
ist, dass Elemente, die zwei äussere und eine mittlere Beobachtung darstellen, auch von der ganzen Reihe wenig abweichen können: so würde folglich der Kegelschnitt fast ganz unbestimmt sein...“

Wie richtig Dr.\ \textsc{Gauss} hierin geurtheilt habe, liegt nun bei Wiederauffindung der \textit{Ceres} am Tage, und die nahe Übereinkunft seiner Elemente mit den jetzigen Beobachtungen bestätigt es. Ja, ohne die scharfsinnigen Bemühungen und Berechnungen des Dr.\ \textsc{Gauss} hätten wir die \textit{Ceres} vielleicht nicht wieder gefunden, der grössere und schönere Theil des Verdienstes gebührt daher ihm. Derselben Meinung ist auch der würdige Dr.\ \textsc{Olbers}; er schreibt an den Herausgeber: „Mit Vergnügen werden Sie bemerkt haben, wie genau Dr.\ \textsc{Gauss}'s Ellipse mit den Beobachtungen der \textit{Ceres} stimmt. Melden Sie doch dies diesem würdigen Gelehrten, unter Bezeugung meiner ganz besondern Hochachtung. Ohne seine mühsamen Untersuchungen über die elliptischen Elemente dieses Planeten würden wir diesen vielleicht gar nicht wiedergefunden haben. Ich wenigstens hätte ihn nicht so weit ostwärts gesucht.“
\begin{flushright}[\textsc{von Zach}.]\end{flushright}

\bigskip\Needspace{10\baselineskip}
\pdfbookmark[1]{Fortgesetzte Nachrichten. Ceres Ferdinandea. 1802 März}{ceres180203}
\begin{center}
{\small FORTGESETZTE NACHRICHTEN ÜBER DEN NEUEN HAUPT-PLANETEN.\par}
CERES FERDINANDEA.\par
\medskip\hrule\medskip
Monatliche Correspondenz. Band V. S.\ 263—282. 1802 März.\par
\medskip\hrule
\end{center}\medskip
—\quad —\quad —\quad „Ich kann nicht umhin, schreibt dieser verdienstvolle Gelehrte [\textsc{Gauss}], hier noch zu erwähnen, was für eine Wohlthat für die Astronomie bei dieser Gelegenheit das Dasein einer Zeitschrift, wie die M.\ C.\ gewesen ist. Mit welcher Lauigkeit und Gleichgültigkeit würde man nicht \textsc{Piazzi}'s Entdeckung aufgenommen haben, wenn Sie nicht durch Ihre Zeitschrift alle Nachrichten darüber gesammelt, auf das schnellste verbreitet, das allgemeine Interesse erweckt, Gründe und Gegengründe abgewogen, und den Planetismus dieses Gestirns zur höchsten Wahrscheinlichkeit gebracht hätten. Wahrscheinlich hätten nur wenige Astronomen sich die Mühe gegeben, es wieder aufzusuchen, da selbst aller jetzigen Astronomen Lehrer und Meister noch vor kurzem den neuen Planeten so stark bezweifelte?“ —\quad —\quad —

—\quad —\quad Wir haben unsern Lesern in dem December-Hefte vorigen Jahres S.\ 646 [S.\ 203 d.\ B.] alle nach und nach von Dr.\ \textsc{Gauss} berechnete und bis zum \textit{vierten mal} nach \textsc{Piazzi}'schen Beobachtungen verbesserte Elemente der Bahn mitgetheilt. In diesen \textit{vierten} Elementen hatte Dr.\ \textsc{Gauss} in der Folge noch folgende kleine Veränderungen gemacht:
\begin{quote}
$\AscendingNode\ 81^0\ 0'\ 13''39$, Sonnenferne $326^0\ 27'\ 41''93$, Epoche $77^0\ 36'\ 38''35$, Neigung $10^0\ 37'\ 1''98$.
\end{quote}
Das übrige blieb ungeändert. —\quad —

—\quad —\quad —\quad Wir haben es schon in dem vorhergehenden Februar-Hefte S.\ 178 erwähnt, dass uns Dr.\ \textsc{Gauss}, noch vor der Auffindung der \textit{Ceres}, eine zum \textit{fünften mal} verbesserte Bahn dieses Planeten zugeschickt hatte. Wir wollen daher auch diese vorerst zur historischen Kenntniss unserer astronomischen Leser gelangen lassen. Dr.\ \textsc{Gauss} unterwarf die \textsc{Piazzi}'schen Beobachtungen einer wiederholten sorgfältigen Berechnung der elliptischen Bahn, wobei er auf Nutation, Aberration und Parallaxe auf das genaueste Rücksicht nahm. Hieraus entstanden folgende fünfte Elemente:
\SourceLeaf{211}%
\begin{center}
(V)\par\smallskip
{\small\setlength{\tabcolsep}{2pt}
\begin{tabular}{l r}
Sonnenferne $\ .\ .\ .\ .\ .\ .\ .\ .\ .$&$324^0\ 37'\ 11''$\\
$\AscendingNode\ .\ .\ .\ .\ .\ .\ .\ .\ .\ .\ .\ .$&$80\ 59\ 12$\\
Neigung $\ .\ .\ .\ .\ .\ .\ .\ .\ .\ .$&$10\ 37\ 9''55$\\
Logarithmus der halben grossen Axe&$0.4446804$\\
Excentricität $\ .\ .\ .\ .\ .\ .\ .\ .\ .$&$0.0879111$\\
tägliche mittl.\ helioc.\ mittl.\ Bewegung&$763''950$\\
Epoche 1800 31. December $\ .\ .\ .\ .$&$78^0\ 5'\ 16''6$
\end{tabular}}
\end{center}
Es wäre überflüssig, hier die schöne Übereinstimmung dieser Elemente mit den \textsc{Piazzi}'schen Beobachtungen zu zeigen, welche Dr.\ \textsc{Gauss} auf das sorgfältigste berechnet und bis auf ein Paar Secunden genau dargestellt hatte.

Als ich dem Dr.\ \textsc{Gauss} die Nachricht von der glücklichen Auffindung der so sehnlichst erwarteten \textit{Ceres}, und meine drei ersten Beobachtungen derselben mitgetheilt hatte, so war das erste, was er nach Empfang derselben that, dass er sie sogleich nach seinen oben angezeigten V\textsuperscript{ten} Elementen berechnete. Er fand den Fehler bei der ersten Beobachtung vom 7. December in AR $+24'\ 8''$; bei der zweiten vom 11. Januar $+30'\ 53''$; bei der dritten vom 16. Januar $+31'\ 53''$. Nach seinen IV\textsuperscript{ten} (verbesserten) Elementen weichen sie nach einem Überschlage in folgender Ordnung ab: $+14\frac12$ Minuten $+19\frac34$ Minuten $+20\frac12$ Minuten. Dass diese IV Elemente der Wahrheit etwas näher kommen, als die V\textsuperscript{ten}, hält Dr.\ \textsc{Gauss} für Zufall; vielleicht ist es aber auch zum Theil Folge der Einwirkung der Planeten-Störungen bei den \textsc{Piazzi}'schen Beobachtungen, besonders auf die Breiten. Soviel Dr.\ \textsc{Gauss} aus den zwei ersten \textsc{Olbers}'schen Beobachtungen schliessen konnte, welche er inzwischen erhalten hatte, so stimmten die Breiten mit beiden Elementen bis auf ein Paar Minuten, und die Differenz in den Declinationen war hauptsächlich nur Folge von den Längen-Differenzen. Indessen ehe Dr.\ \textsc{Gauss} weitere Beobachtungen von mir erhielt, konnte er der Versuchung nicht widerstehen, eine provisorische Correction seiner Elemente vorzunehmen. Er hielt sich dabei an meine beiden geraden Aufsteigungen vom 7. December und 16. Januar, und so fand er folgende VI\textsuperscript{ten} Elemente, nach einer mit aller möglichen Sorgfalt geführten Rechnung, welche um so nöthiger war, da der geringste Fehler die ganze Arbeit unnütz gemacht hätte.
\begin{center}
{\small\setlength{\tabcolsep}{2pt}
\begin{tabular}{l l}
Epoche 1801 Palermer Meridian $\ .\ .$&$77^0\ 24'\ 55''9$\\
—\quad1802\quad—\quad— $\ .\ .$&$155\ 33\ 35.1$\\
Tägliche mittlere tropische Bewegung&$770''7376$\\
Tropische Umlaufszeit $\ .\ .\ .\ .\ .\ .$&1681 Tage 12 St.\ 9 Min.\\
Logarithmus der halben grossen Axe&$0.4421189$
\end{tabular}
\par\smallskip
$\left.\begin{array}{lr}
\text{Sonnenferne}\ .\ .\ .\ .\ .\ .\ .\ .\ .&326^0\ 14'\ 45''\\
\AscendingNode\ .\ .\ .\ .\ .\ .\ .\ .\ .\ .\ .\ .&80\ 58\ 55
\end{array}\right\}\quad\begin{gathered}
\text{31. December 1800,}\\
\text{siderisch ruhend angenommen}
\end{gathered}$
\par\smallskip
\begin{tabular}{l r}
Excentricität $\ .\ .\ .\ .\ .\ .\ .\ .\ .$&$0.08086253$\\
grösste Mittelpunktsgleichung $\ .\ .\ .$&$9^0\ 16'\ 23''$\\
Neigung $\ .\ .\ .\ .\ .\ .\ .\ .\ .\ .\ .$&$10\ 37\ 51$
\end{tabular}}
\end{center}
Es lag in der Natur der bei dieser Rechnung gebrauchten Methode, dass durch diese Elemente die \textsc{Piazzi}'schen Beobachtungen fast eben so scharf dargestellt werden, als durch die V\textsuperscript{ten} Elemente. Allein es war nicht zu erwarten, dass gleich durch diese erste Verbesserung die Fehler bei den beiden neuen Beobachtungen sogleich von einem halben Grad auf Null gebracht würden. Inzwischen als Dr.\ \textsc{Gauss} nach diesen neuen Elementen die beiden dabei zum Grunde gelegten geraden Aufsteigungen berechnet hatte, fand er doch die Übereinstimmung grösser, als er gehofft hatte. Der Fehler bei der ersten Beobachtung war nur $+3''0$, in der letzten $+9''9$. Dr.\ \textsc{Gauss} hatte nun ferner drei neue Beobachtungen von mir, welche ich ihm unterdessen zugeschickt hatte, nach diesen Elementen berechnet, und fand folgende Übereinstimmung der geraden Aufsteigungen.
